\scanpage{8}
% L01–02
$\mathbb E^n$ sur $\mathbb F_p$, $\Gamma_n=\operatorname{Gl}(n,\mathbb F_p)_{\mathbb F_p}\subset\operatorname{Gl}(n)_{\mathbb F_p}$ \inserted{opère sur $\mathbb E^n$}
% L03
\[ \unsure{Q|\mathbb Q}_n=\mathbb E^n/\Gamma_n=\operatorname{Spec}(A_n)\qquad A_n=\mathbb F_p[X_1,\ldots,X_n]^{\Gamma_n} \]
% L04–05
Soit $U_n\subset\mathbb E^n$ l'ouvert formé des ``syst. lin. indép. sur $\mathbb F_p$'' i.e. l'\deleted{\unsure{ens}} \unsure{plus} grand ouvert de $\mathbb E^n$ sur lequel $\Gamma_n$ opère librement, $Q_n^0=U_n/\Gamma_n$

% L06
\uline{Lemme} \deleted{On a}
% L06
{\scriptsize
\[
\deleted{\det\begin{pmatrix}X_1&X_2&\cdots&X_n\\X_1^p&X_2^p&\cdots&X_n^p\\\ldots&&&\\X_1^{p^{n-1}}&X_2^{p^{n-1}}&\cdots&X_n^{p^{n-1}}\end{pmatrix}=
\prod_{\alpha\in\mathbb F_p^{n-1}}(X_n+\sum_1^{n-1}\alpha_iX_{\unsure{n|i}})
\prod_{\alpha\in\mathbb F_p^{\unsure{n-2|n-1}}}X_{n-2}+\sum_1^{\unsure{n-2}}\alpha_iX}
\]}
% L07
On a dans $\mathbb F_p[X_1,\text{--},X_n]$
% L08–11
{\scriptsize
\[
\begin{aligned}
\underset{\substack{\shortparallel\,\mathrm{d\acute ef}\\ \Delta_n(X_1,\text{--},X_n)}}{\det\begin{pmatrix}
X_1&\cdots&X_n\\X_1^p&\cdots&X_n^p\\\ldots&&\\X_1^{p^{n-1}}&\cdots&X_n^{p^{n-1}}\end{pmatrix}}
&=\prod_{(\alpha_1,\ldots,\alpha_{n-1})\in\mathbb F_p^{n-1}}\left(X_n+\sum_1^{n-1}\alpha_i X_i\right)
\prod_{(\alpha_1,\ldots,\alpha_{n-2})\in\mathbb F_p^{n-2}}\left(X_{n-1}+\sum_1^{n-2}\alpha_iX_i\right)\ldots\\
&\qquad\ldots\left(\prod_{\alpha_1\in\mathbb F_p^1}(X_2+\alpha_1X_1)\right)X_1\\
&=\prod_{0\leq j\leq n-1}\deleted{\unsure{X_{j+1}}}\prod_{\alpha\in k^j}\left(X_{j+1}+\sum_{i=1}^j\alpha_i X_i\right)
\end{aligned}
\]}

% L12–16
\uline{NB} À un facteur $\in\mathbb F_p^*$ près, le deuxième mb est le produit des combinaisons linéaires non \deleted{\unsure{toutes}} id. nulles des $X_i$ à coeff dans $\mathbb F_p$, \deleted{\unsure{en}} \deleted{\unsure{ne pre}} prenant ces combinaisons mo\uline{dulo} \deleted{\unsure{homothétie}} \inserted{multiplication} par un $c\in\mathbb F_p^*$. C'est un polynôme \inserted{$P$} de degré
% L17
\[ 1+p\ \text{---}+p^{n-1}=\frac{p^n-1}{p-1}=\operatorname{card}P^{n}(\mathbb F_p) \]
% L18–23
donc de degré égal au premier, $Q$. Or \deleted{\unsure{les}} l'ens. des zéros de \deleted{produits} \inserted{chaque facteur $P_\lambda$ de $P$} est contenu dans celui de $Q$, \deleted{\unsure{et comme $P$}} donc \inserted{($P_\lambda$ étant premier)} $P_\lambda$ divise $Q$, donc (les $P_\lambda$ étant premiers 2 à 2 \deleted{\unsure{dist}} non associés) $Q$ est divisible par $P=\prod_\lambda P_\lambda$, donc, étant de $=$ degré, est de la forme $cP$, $c\in\mathbb F_p^*$. Reste à déterminer $c$ \ldots
% archive sheet number: 7
